\section{边值问题：从打靶到全局离散}
\label{sec:08-Numerical-Analysis/06-ODE/04}

一根金属棒，两端分别浸在冰水和沸水里，中间有一段内部热源。要的是稳态时棒上的温度分布。方程写出来毫无难度，

\[
-y''(x)=f(x),\qquad y(0)=\alpha,\quad y(1)=\beta,
\]

可前三节的所有方法在它面前都没法启动。左端知道 \(y(0)\)，不知道 \(y'(0)\)；缺了初始斜率，第一步就迈不出去。

约束不在一端，而是分散在两端。这个改动看着不大，它取消的却是整套框架的前提：未知量不再是``下一步的状态''，而是整条曲线。

\subsection{打靶法把缺失的初值变成一个求根问题}

既然只缺一个数，那就把它当未知量。设 \(y'(0)=s\)，从 \((\alpha,s)\) 出发解一个标准的初值问题，一直积到 \(x=1\)，记落点为 \(y(1;s)\)。定义

\[
F(s)=y(1;s)-\beta,
\]

问题就化成了解 \(F(s)=0\)。每一次 \(F\) 的求值是跑一遍完整的 IVP，外层用二分、割线或 Newton 迭代——正是\hyperref[sec:08-Numerical-Analysis/02-Root-Finding/02]{``不动点与 Newton 法''}里那套工具，只不过被求值的函数昂贵得多。

\begin{center}
\begin{tikzpicture}[x=5.8cm,y=6cm,font=\small]
  \draw[->,black!55] (-0.04,0) -- (1.2,0) node[right,font=\footnotesize] {\(x\)};
  \draw[black!30,dashed] (1,-0.36) -- (1,0.38);
  \draw[thick,black!45] plot[domain=0:1,samples=60] (\x,{0.2*\x-\x*\x/2});
  \draw[thick,black!45] plot[domain=0:1,samples=60] (\x,{0.8*\x-\x*\x/2});
  \draw[very thick,blue!70] plot[domain=0:1,samples=60] (\x,{0.5*\x-\x*\x/2});
  \fill[black!70] (0,0) circle (1.8pt);
  \node[left,font=\footnotesize] at (-0.02,0) {\(\alpha\)};
  \draw[red!70,line width=1.4pt] (1,0) circle (2.4pt);
  \node[right,font=\footnotesize,red!75] at (1.04,0.09) {\(\beta\)：必须命中};
  \fill[black!55] (1,-0.3) circle (1.6pt);
  \fill[black!55] (1,0.3) circle (1.6pt);
  \node[right,font=\footnotesize,black!60] at (1.04,-0.3) {\(s\) 太小，打低了};
  \node[right,font=\footnotesize,black!60] at (1.04,0.3) {\(s\) 太大，打高了};
  \node[above,font=\footnotesize,blue!70] at (0.5,0.14) {\(s=s^{\star}\)};
  \draw[->,>=stealth,black!45] (0.06,-0.115) -- (0.012,-0.012);
  \node[right,font=\footnotesize,black!65] at (0.05,-0.145) {同一起点，三种不同的初始斜率};
\end{tikzpicture}

\emph{打靶法把两端问题重新装回单向推进的框架：从左端以不同斜率射出，读右端的偏差，交给求根法去调那一个数。}
\end{center}

这条路线的吸引力在于复用。成熟的 IVP 求解器、自适应步长、事件处理、刚性方法，全都不用改一行就能用上，外面只需套一层求根循环。低维、稳定、对初值不过分敏感的问题上，它是最省事的做法。

麻烦藏在``对初值敏感''这几个字里。若方程含有增长模态，特征指数为 \(\mu>0\)，那么初始斜率的扰动 \(\delta\) 传到右端会变成大约 \(\delta e^{\mu(b-a)}\)。区间稍长一点，\(F(s)\) 的斜率就大到荒唐：右端要求的精度是 \(10^{-8}\) 时，\(s\) 需要被确定到 \(10^{-8}e^{-\mu(b-a)}\)，很快低于机器精度，射出去的每一条轨迹在数值上都无法区分。这时 \(F\) 的病态不是求根法的锅，是这个转化本身把问题的条件数放大了 \(e^{\mu(b-a)}\) 倍。

补救方向是别让任何一段积得太长。multiple shooting 把区间切成若干小段，每段引入自己的初值作为未知量，段与段之间用连续性条件联立，最后解一个中等规模的非线性方程组。指数放大因子从 \(e^{\mu(b-a)}\) 降到每一小段的 \(e^{\mu\Delta}\)，代价是未知量从一个变成几十个。它介于打靶与下一小节之间。

\subsection{全局离散把整条曲线一次解出来}

另一条路线干脆放弃推进。取网格 \(x_i=ih\)，把二阶导数换成中心差分，

\[
y''(x_i)\approx\frac{y_{i-1}-2y_i+y_{i+1}}{h^{2}},
\]

它由 \(y(x_i\pm h)\) 的两个 Taylor 展开相加得到，奇数阶项相消，余项是 \(\frac{h^{2}}{12}y^{(4)}(\xi_i)\)，所以是二阶近似。代入方程，每个内点给出一个方程

\[
-y_{i-1}+2y_i-y_{i+1}=h^{2}f(x_i),
\]

连同两端的边界值，凑成一个三对角线性系统 \(A_h y=b\)。所有节点值同时被解出来，没有先后。

误差分析走的仍是``一致性加稳定性推出收敛''那条老路，但要点值得写清楚。把精确解的采样代入离散方程，残差 \(\tau=O(h^{2})\)；误差 \(e\) 满足 \(A_h e=-\tau\)，于是

\[
\norm{e}\le\norm{A_h^{-1}}\,\norm{\tau}=O(h^{2}).
\]

对账：这里的 \(A_h\) 是已经除过 \(h^{2}\) 的那一版，即对算子 \(-d^2/dx^2\) 的离散；关键条件是 \(\norm{A_h^{-1}}\) 随 \(h\to0\) 一致有界，这条与连续问题的适定性对应，光有 Taylor 展开推不出收敛。范数取任一向量范数及与之相容的矩阵范数。

对照打靶法，差别一目了然。全局离散不沿任何方向传播误差，不存在指数放大那回事；换来的是要解一个 \(N\) 元方程组。一维三对角系统用 Thomas 消元是 \(O(N)\)，几乎免费；推广到二维三维，或者非线性 BVP 每次 Newton 迭代都要解一次，矩阵规模就成了主要成本，这时线性求解器的选择重新变成瓶颈，回到\hyperref[sec:08-Numerical-Analysis/05-Numerical-Linear-Algebra/04]{``大型稀疏系统与 Krylov 迭代''}的地盘。

两条路线的取舍可以一句话说完：打靶用小规模的求根换来了对初值的指数敏感，全局离散用大规模的线性代数买断了这份敏感。问题维数低、区间短、动力学温和时选前者；区间长、含增长模态、或者本来就要在多维区域上求解时选后者。

配置法与有限元是全局离散的推广。前者先选一族有限维基函数（分段多项式、样条、全局谱多项式），要求方程在若干配置点上精确成立；后者用弱形式控制残差。光滑问题上谱配置的收敛快得惊人，几何不规则或解有局部剧变时，有限元与自适应网格更合适。选择取决于解的光滑性、维数和边界形状，不取决于阶数高低。

\subsection{边界条件能让问题无解，也能让它有无穷多解}

有一类失败与离散无关，出在连续问题上。看纯 Neumann 问题

\[
-y''=f,\qquad y'(0)=y'(1)=0 .
\]

把方程在 \([0,1]\) 上积分，左边得到 \(-\bigl(y'(1)-y'(0)\bigr)=0\)，于是必须有

\[
\int_0^1 f(x)\,dx=0 .
\]

这条兼容条件不满足，方程就没有解；满足了，解也只确定到一个相加常数——若 \(y\) 是解，\(y+C\) 同样满足方程和两个边界条件。

离散之后这件事会以另一副面孔出现：矩阵奇异，求解器报告条件数是无穷。此时在对角线上加一个小量，程序确实不再报错，返回的却是另一个问题的解。正确的处理是加一条均值归一化，或者固定某一个自由度，并且回头检查 \(f\) 的兼容性。数值上的奇异是连续问题在告诉你它本来就欠定，不是实现的 bug。

非线性 BVP 还可能有多个解——梁的屈曲、燃烧的多稳态、反应扩散的分支，都是这样。Newton 从哪里出发，决定收敛到哪一支。代数残差降到 \(10^{-12}\) 只说明离散系统解得干净，它对``这一支是不是物理上要的那一支''不发表任何意见。

\subsection{三层验证缺一不可}

BVP 的调试之所以容易走偏，是因为可查的东西分在三个层次上，而它们互相不能替代。代数层查离散方程的残差，回答``方程组解对了吗''。边界层查每个边界条件是否按正确的尺度被满足，Neumann 条件尤其容易在离散时掉一个 \(h\) 的因子而不报错。连续层靠网格加密：把 \(h\) 减半，误差应当按预期的阶下降，守恒量与对称性应当保持。

只做前两层，可能把一个粗糙但错误的离散问题解得非常精确。这类错误不会抛异常，也不会让残差变大，只会让曲线安静地偏在一边。

至此，从初值到边值、从显式到隐式的路径走完了一遍。贯穿其中的是同一件事：局部信息足以定义一步，但一步之外的行为由累积、放大和约束的位置共同决定，而这三者都不写在方程的右端项里。
